\chapter{Energies and the ASE calculator}
\label{ch:energies}

Once $\rho$ and $\tau$ have been predicted, the \pyclass{poraque.physics}
module turns them into energies, and the \pyclass{Poraque} calculator wraps the
whole chain behind the ASE calculator protocol.

\begin{pwarn}
Read Section~\ref{sec:not-a-dft-energy} before comparing any of these numbers
to a DFT result. They are \textbf{pseudo-valence} energies, and on the current
models the error on energy \emph{differences} is larger than the differences
themselves.
\end{pwarn}

\section{The energy expression}

\[
  E = \underbrace{\int\tau\,d^3r}_{\Ts}
    + \underbrace{\int\rho \vext\,d^3r + E_{\alpha Z}}_{E_{\rm ext}}
    + \underbrace{\tfrac12\int\rho \vH\,d^3r}_{E_{\rm H}}
    + E_{\rm xc}[\rho]
    + E_{\rm Ewald}
\]

\begin{center}
\begin{tabular}{@{}lp{8.4cm}@{}}
\toprule
Term & How it is obtained \\
\midrule
$\Ts$ & integral of the predicted \file{TAUCAR} \\
$E_{\rm ext}$ & $\int\rho\vext$, with $\GG = 0$ excluded \\
$E_{\alpha Z}$ & the finite $\GG = 0$ remainder, from the \file{POTCAR}
\opt{PSCORE} \\
$E_{\rm H}$ & Poisson solve in reciprocal space, $\vH(\GG{=}0) = 0$ \\
$E_{\rm xc}$ & PBE by default; LDA available (Section~\ref{sec:xc}) \\
$E_{\rm Ewald}$ & Ewald sum over the pseudo-ions in a neutralizing
background \\
\bottomrule
\end{tabular}
\end{center}

\subsection{The \texorpdfstring{$\GG = 0$}{G = 0} bookkeeping}
\label{sec:g0}

Each of the three electrostatic terms diverges at $\GG = 0$, and the
divergences cancel in a neutral cell. Poraqu\^e follows the standard plane-wave
treatment: drop $\GG = 0$ from all three, then add back the finite remainder of
the electron-ion term,
\[
  E_{\alpha Z} = \frac{N_{\rm elec}}{\Omega}\sum_s N_s\,\mathrm{PSCORE}_s ,
\]
which is VASP's \opt{alpha Z}. It is of order eV \emph{per atom}, so it is read
from the \file{POTCAR} tables when they are available and reported as
\opt{None} when they are not --- never silently assumed to be zero.

\begin{pnote}
\opt{components.missing} lists any term that was skipped, and printing the
components emits \opt{incomplete:} when that list is non-empty. A total that is
missing a term says so rather than looking complete.
\end{pnote}

\section{Choosing the exchange-correlation functional}
\label{sec:xc}

\begin{center}
\begin{tabular}{@{}lll@{}}
\toprule
\opt{functional} & Exchange & Correlation \\
\midrule
\opt{"pbe"} (\textbf{default}) & PBE & PBE \\
\opt{"lda"} & Dirac & PW92 \\
\opt{"pbe-x"} & PBE & --- \\
\opt{"lda-x"} / \opt{"x-only"} & Dirac & --- \\
\opt{"none"} & --- & --- \\
\bottomrule
\end{tabular}
\end{center}

PBE is the default because it matches the reference data: the calculations
Poraqu\^e ingests use \opt{PAW\_PBE} pseudopotentials with \opt{LEXCH = PE}, so
$\rho$ and $\tau$ are PBE quantities. Evaluating an LDA $E_{\rm xc}$ on a PBE
density is neither a PBE energy nor an LDA one. On the reference Au supercells
the two differ by $-0.92$ eV/atom ($0.65\,\%$ of $E_{\rm xc}$), so the choice
is not cosmetic.

Both are validated against the limit that defines them: on a uniform density
PBE reduces to LDA \emph{bit-exactly}, since $F_{\rm x}(0) = 1$ and $H \to 0$.
Dirac exchange reproduces $-0.458165293/r_s$ Ha per electron exactly, and PW92
matches the published table to $10^{-6}$.

\begin{pwarn}
PBE is semilocal and needs $\nabla\rho$. On a \emph{predicted} field that
gradient carries the network's noise and the enhancement factor amplifies it;
on a band-limited grid it also aliases wherever the density has sharp core
peaks. The extra physics is only worth having once the density is good enough
for its gradient to mean something --- which, on the current models, it is not
(Section~\ref{sec:not-a-dft-energy}).
\end{pwarn}

Set it wherever the energy is computed:

\begin{lstlisting}
EnergyCalculator.from_potential(vext, functional="lda")
Poraque("ext2chg.pfno", "chg2tau.pfno", potcar_dir=..., functional="lda")
\end{lstlisting}
\begin{lstlisting}
poraque-inference run/ \
    --functional lda
\end{lstlisting}

\section{Computing energies from files}

\begin{lstlisting}
from poraque.fields import ChargeDensity, ExternalPotential, \
                           FieldGrid, KineticEnergyDensity
from poraque.fields.vasp.potcar import Potcar
from poraque.physics import EnergyCalculator

grid = FieldGrid.from_file("run/CHGCAR")
rho  = ChargeDensity.read("run/CHGCAR", grid=grid)
tau  = KineticEnergyDensity.read("run/TAUCAR", grid=grid)
vext = ExternalPotential.from_calculation("run", grid=grid)

potcar = Potcar.from_file("run/POTCAR", parse_tables=True)
calc = EnergyCalculator.from_potential(
    vext, pscore={e.element: e.pscore for e in potcar})

print(calc.compute(rho, tau, vext))
\end{lstlisting}

\begin{lstlisting}
  kinetic      T_s            6207.068300 eV
  external     E_ext        -12634.744762 eV
  alpha Z                     1752.512988 eV
  Hartree      E_H            3230.363304 eV
  xc (pbe)                   -3845.824650 eV
  Ewald                     -25949.345827 eV
  ----------------------------------------
  potential                 -37447.038948 eV
  TOTAL                     -31239.970647 eV
  electrons                    297.000001
\end{lstlisting}

\begin{ptip}
\opt{electrons} is the cheapest available diagnostic. It should equal the total
\opt{ZVAL} --- 297 for the 27 Au atoms above --- to a few parts in $10^{4}$. A
predicted density that has drifted off it invalidates every electrostatic term,
since all of them are linear or quadratic in $\rho$.
\end{ptip}

Plain arrays are accepted as readily as field objects, so a prediction can be
scored without being written to disk first.

\section{The ASE calculator}

\pyclass{Poraque} behaves like MACE or NequIP at the interface:

\begin{lstlisting}
from ase.build import bulk
from poraque.calculator import Poraque

atoms = bulk("Au", "fcc", a=4.08, cubic=True)
atoms.calc = Poraque("models/poraque_models.pfno", potcar="POTCAR")

energy = atoms.get_potential_energy()
print(atoms.calc.components)          # the full decomposition
rho = atoms.calc.fields["density"]    # the predicted CHGCAR
\end{lstlisting}

Each call runs \opt{Atoms} $\to$ \opt{Structure} $\to$ \opt{FieldGrid} $\to$
$\vext$ $\to$ $\rho$ $\to$ $\tau$ $\to$ $E$. The three fields stay on the
calculator afterwards, so a prediction can be written out with
\opt{rho.write("CHGCAR")}.

\begin{center}
\begin{tabular}{@{}lp{8.0cm}@{}}
\toprule
Argument & Meaning \\
\midrule
\opt{ext2chg}, \opt{chg2tau} & Checkpoint paths, or loaded
\pyclass{FieldOperator}s. \\
\opt{potcar} & A single \file{POTCAR} covering the species present. \\
\opt{potcar\_dir} & A \file{POTCAR} \emph{library}, searched per
composition. \\
\opt{charges} & \opt{\{element: Z\_val\}}, used only without a
\file{POTCAR}. \\
\opt{resolution} & Longest grid axis; defaults to the checkpoint's training
resolution. \\
\opt{functional} & Exchange-correlation approximation, default \opt{"pbe"}. \\
\opt{device} & \opt{auto}, \opt{cuda}, \opt{mps} or \opt{cpu}. \\
\bottomrule
\end{tabular}
\end{center}

\subsection{Supplying pseudopotentials}

\opt{potcar} is right when the composition is fixed. For a calculator that must
serve \textbf{arbitrary} structures, point \opt{potcar\_dir} at a \file{POTCAR}
library: the entries for whatever elements an \opt{Atoms} contains are
assembled on demand and cached per composition.

\begin{lstlisting}
atoms.calc = Poraque("models/poraque_models.pfno",
                     potcar_dir="/opt/vasp/potpaw_PBE")
\end{lstlisting}

Recognised layouts, in preference order --- each accepting a \file{.gz} or
\file{.Z} suffix:

\begin{enumerate}
  \item \file{<dir>/Au/POTCAR} --- what VASP ships;
  \item \file{<dir>/Au\_pv/POTCAR} --- a variant, used only if it is the
    \emph{only} candidate, with a warning;
  \item \file{<dir>/POTCAR.Au} or \file{<dir>/Au.POTCAR} --- flat.
\end{enumerate}

\begin{pnote}
If several variants match --- \opt{Fe\_pv} and \opt{Fe\_sv}, say --- the lookup
\textbf{raises} rather than picking one. They differ in \opt{ZVAL} and
therefore in every energy, so the choice is the user's; name the one you want
with an explicit \opt{potcar}.
\end{pnote}

\begin{pwarn}
With no \file{POTCAR} at all the external potential falls back to the
\emph{Gaussian} pseudo-ion model, which differs from the tabulated potential by
a relative $L^{2}$ of about $0.13$ --- against $2\times10^{-5}$ for the
tabulated one. The operators were trained on tabulated potentials, so the
Gaussian model feeds them an input outside their training distribution. The
calculator warns once and proceeds; treat the result as a smoke test, not a
prediction.
\end{pwarn}

\subsection{Forces and stress}

\opt{get\_forces()} returns the analytic \textbf{Hellmann--Feynman} force: the
electron--ion term $-\int\rho\,\partial\vext/\partial\Rv$ evaluated in
reciprocal space from the same \file{POTCAR} form factor the potential was
built from, plus the analytic Ewald force. Both are differentiated in closed
form rather than by finite differences --- on a fixed grid those would be
dominated by the grid's own discontinuity as atoms cross voxel boundaries.

The implementation is exact: each term agrees with a central finite difference
of the energy it differentiates to $10^{-7}$ eV/\AA.

\begin{pwarn}
The \emph{physics} is incomplete for PAW datasets. Hellmann--Feynman is the
whole force only for a \emph{local} pseudopotential; a PAW calculation adds a
projector force and a one-centre force, and neither is recoverable from $\rho$,
$\tau$ and $V_{\rm loc}$ on a grid. Measured against VASP on gold, using
VASP's own density: MAE $0.83$ eV/\AA\ against forces of $1.66$ eV/\AA. The
magnitude is right, the direction is not, so \textbf{geometry optimisation and
molecular dynamics remain unavailable}.

The cancellation is why it is delicate: the electron--ion and ion--ion terms
are each $\approx 100$ eV/\AA\ and cancel to $\approx 0.5$ eV/\AA, so a
relative error in $\rho$ arrives in the force amplified roughly $200$-fold.
\end{pwarn}

\opt{get\_stress()} still raises: the stress needs the energy's response to a
strain, which deforms the cell \emph{and} the grid the fields live on.

\opt{implemented\_properties} is
\opt{["energy", "free\_energy", "forces"]}. The first two are the same number:
no electronic smearing enters this pipeline, and ASE optimizers request
\opt{free\_energy} by name.

\subsection{Cohesive energies}
\label{sec:cohesive}

A total energy means nothing on its own. It is referenced to whatever the
pseudopotential generator chose, and Poraqu\^e's totals additionally carry an
offset of order $10^{3}$ eV per atom from the absent PAW one-centre terms. The
\textbf{cohesive energy} removes the arbitrary part,

\begin{equation*}
\Delta E = E_{\rm total} - E_{\rm ref},
\qquad
E_{\rm ref} = \sum_i E_{\rm iso}(Z_i),
\end{equation*}

with $E_{\rm iso}$ the energy of one isolated atom of that species. What
survives is the energy released on assembling the solid from free atoms --- the
bonding, and nothing else.

Supply a directory holding one subdirectory per element, each an ordinary
single-point calculation of one atom in a large box:

\begin{lstlisting}
data/vasp/ref/
    Au/     POSCAR POTCAR OSZICAR OUTCAR CHGCAR TAUCAR
    N/      ...
\end{lstlisting}

\begin{lstlisting}[language=Python]
atoms.calc = Poraque("models/poraque_models.pfno",
                     potcar="POTCAR",
                     references="data/vasp/ref")

atoms.get_potential_energy()                   # total, as always
atoms.calc.get_cohesive_energy()               # dE, in eV
atoms.calc.get_cohesive_energy(per_atom=True)  # eV/atom
\end{lstlisting}

\subsubsection{Which reference to subtract}

\pyclass{ReferenceEnergies.from\_directory} takes a \opt{method}, and the
choice matters more than anything else here.

\begin{center}
\begin{tabular}{@{}lll@{}}
\toprule
\opt{method} & $E_{\rm iso}$ from & Au cohesive energy \\
\midrule
\opt{"poraque"} (default) & Poraqu\^e's own expression & $\mathbf{-1.9}$ eV/atom \\
\opt{"code"} & VASP \file{OUTCAR}/\file{OSZICAR} & $-1157$ eV/atom \\
\bottomrule
\end{tabular}
\end{center}

A cohesive energy is meaningful only when the two energies being subtracted
carry the \emph{same} systematic error. Subtracting VASP's atomic energy from
Poraqu\^e's total leaves the PAW offset entirely intact; subtracting
Poraqu\^e's own atomic energy cancels it, because the same terms are missing
from both sides. Hence the default. Use \opt{"code"} when the reference
calculations are what you want to compare \emph{against}.

\begin{pnote}
Referencing changes neither the forces nor energy differences at fixed
composition, and both statements are exact identities rather than
approximations. $E_{\rm ref}$ depends on composition and not on coordinates, so
$\nabla_{\Rv} E_{\rm ref} = 0$; and two structures with the same formula share
$E_{\rm ref}$, which cancels bit-for-bit in $\Delta E_1 - \Delta E_2$. Its
value is in comparisons \emph{across} compositions --- binding energies,
cohesive energies per atom, formation energies --- where the offset does not
cancel and without which the comparison is undefined.
\end{pnote}

\subsection{The Hartree potential}
\label{sec:hartree-field}

$\vH$ is \textbf{solved, not predicted}. Poisson's equation relates it to
$\rho$ exactly, and on a periodic plane-wave grid the reciprocal-space form is
the solution rather than an approximation to it:

\begin{equation*}
\nabla^2 \vH = -4\pi e^2\rho
\quad\Longleftrightarrow\quad
\vH(\GG) = \frac{4\pi e^2\rho(\GG)}{G^2},
\qquad \vH(\GG = 0) = 0 .
\end{equation*}

Two FFTs, exact for any band-limited density. A third learned operator would be
strictly worse: it would inject error into a quantity that has none, and would
not be guaranteed to satisfy the equation that defines it.

\begin{lstlisting}[language=Python]
potential = atoms.calc.get_hartree_potential()
potential.write("LOCPOT")

# v_H + V_ext, the total local potential a plain VASP LOCPOT holds
total = atoms.calc.get_hartree_potential(with_external=True)
\end{lstlisting}

The $\GG = 0$ term is set to zero --- the same neutralising-background
convention $\vext$ uses (\S\ref{sec:g0}), which is what makes the two
potentials addable and what makes $E_{\rm H}$ of a uniform density come out at
exactly zero rather than infinite.

Written to \file{LOCPOT} the field is stored \textbf{unscaled}: unlike
\file{CHGCAR}, which holds $\rho\Omega$, a \file{LOCPOT} holds the potential
itself in eV. From the command line:

\begin{lstlisting}
poraque-inference structure/ --output predictions/ --write-locpot
poraque-inference structure/ --output predictions/ --write-locpot --locpot-total
\end{lstlisting}

The field accessors are symmetric --- \opt{get\_external\_potential()},
\opt{get\_charge\_density()}, \opt{get\_kinetic\_energy\_density()} and
\opt{get\_hartree\_potential()} --- and all four return fields on one shared
grid.

\section{What the number is not}
\label{sec:not-a-dft-energy}

\paragraph{It is not a DFT total energy.} \file{CHGCAR} holds the valence
\emph{pseudo} density and \file{TAUCAR} the valence pseudo kinetic energy
density, so the PAW one-centre terms are absent entirely. For the 27-atom Au
cell above the result is $\approx -31215$ eV against a VASP \opt{TOTEN} of
order $-100$ eV. Nothing in this module can close that gap.

\paragraph{Energy differences are not usable yet either.} Measured on the
current seventeen-structure Au dataset, within each composition group:

\begin{center}
\begin{tabular}{@{}lrr@{}}
\toprule
Source of the fields & MAE on $\Delta E$ & vs.\ signal \\
\midrule
VASP's own $\rho$ and $\tau$ (method ceiling) & 0.18--0.24 eV/atom & 1.5--2$\times$ \\
The shipped operators & 0.87--1.93 eV/atom & 7--15$\times$ \\
\emph{True spread being resolved} & \emph{0.125 eV/atom} & --- \\
\bottomrule
\end{tabular}
\end{center}

The error \textbf{exceeds the signal} in both rows and the correlation with the
true ordering is consistent with zero, so the predicted ranking of structures
carries no information.

The first row is the important one. It is measured with the model removed from
the problem entirely, so it is not a training failure --- it is the neglected
PAW one-centre and non-local energy, which is \emph{not} a per-atom constant.
The electrostatics themselves are correct: $E_{\alpha Z}$, Ewald and
$E_{\rm H}$ reproduce VASP's \opt{PSCENC}, \opt{TEWEN} and \opt{DENC} to better
than $0.01$ eV in a difference, and the Ewald sum reproduces exact Madelung
constants to $10^{-6}$.

The underlying difficulty is scale. The total is a cancellation of terms of
order $10^{4}$ eV down to a result of order $10^{0}$ eV, so resolving $\Delta E$
to $10$ meV/atom needs a relative accuracy of about $10^{-6}$ in the fields.
The operators deliver $10^{-2}$.

\begin{pnote}
The grid is \emph{not} the limiting factor. On reference fields the energy
difference changes by less than $0.2$ eV between \opt{resolution: 32} and the
native $128^{3}$ mesh, and $\Ts$ and the electron count are preserved exactly
--- spectral resampling is an exact band-limited projection.
\end{pnote}

\paragraph{Charge normalisation helps, and is on by default.} The shipped
\opt{ext2chg} operator predicts densities whose electron count drifts by up to
$1.7\,\%$ --- nearly five electrons out of $297$. Every electrostatic term is
at least linear in $\rho$ and $E_{\rm H}$ is quadratic, so that drift alone
moved totals by tens of eV, differently for each structure, and therefore did
not cancel in a difference. Rescaling to the count the pseudopotentials fix
\emph{halved} the end-to-end error on $\Delta E$. It is a repair rather than a
cure: a density whose integral is wrong by $1.7\,\%$ has the wrong shape too,
and rescaling does not fix the shape. Disable it with
\opt{normalize\_density=False} only to diagnose.
